\origpage{98}{ROBERTSON AND SEYMOUR}
\stmt
(8.7) \textit{Let $X$ be an extreme restraint, and let $a, b \in V(\Gamma \cap bd(X))$ be distinct. Let $P$ be a walk in $\Gamma$ from $a$ to $b$, such that no edge or internal vertex of $P$ is in $X$. Then there is a walk of $\Gamma \cap bd(X)$ from $a$ to $b$ which is shorter than $P$.}

\textit{Proof.} We may assume that $P$ is simple. Now by (7.6), $\Sigma - X$ is arc-wise connected. Since $P$ is simple and has no internal vertex in $X$, it follows that $\Sigma - (X \cup U(P))$ has at most two arc-wise connected components, $\Delta_1$ and $\Delta_2$ say, where possibly $\Delta_1 = \Delta_2$. Since only two vertices of $\Gamma \cap bd(X)$ are incident with edges of $P$, and each of them is incident with only one edge of $P$, it follows that there are walks $W_1, W_2$ of $\Gamma \cap bd(X)$ from $a$ to $b$ such that $W_1 + W_2^{-1}$ is a perimeter walk of $X$ and for $i = 1, 2$, $U(\Gamma | W_i) \subseteq \bar{\Delta}_i$.

We suppose for a contradiction that $W_1, W_2$ have length at least that of $P$. Thus $P + W_1^{-1}$ has length at most that of $W_2 + W_1^{-1}$, and hence at most $l(X)$. Let $X_1 = \mathrm{ins}(P + W_1^{-1})$; then $X_1$ is a restraint by (8.5). We claim that $\Delta_1 \subseteq X_1$. For there is a simple walk $W$ from $a$ to $b$ with $\Gamma | W \subseteq \Gamma | W_1$, and hence $P + W^{-1}$ is a simple-closed walk of length $< 2\theta$. It follows that $\mathrm{ins}(P + W^{-1})$ includes one of $\Delta_1, \Delta_2$, and hence $X_1$ includes one of $\Delta_1, \Delta_2$. Suppose that $\Delta_2 \subseteq X_1$. Then $\bar{\Delta}_2 \subseteq X_1$, and therefore $U(\Gamma | W_2) \subseteq X_1$. Hence $bd(X) \subseteq X_1$. By (8.4), $X \cup X_1 \ne \Sigma$, and so $\Sigma - X_1 \nsubseteq X$. Now $\Sigma - X_1$ is arc-wise connected by (8.5) and (7.6), and is disjoint from $bd(X)$; and so $(\Sigma - X_1) \cap X = \emptyset$, that is, $X \subseteq X_1$. But $X \ne X_1$ since $U(P) \subseteq X_1$, and yet by (8.5) $l(X_1) \le l(P + W_1^{-1}) \le l(X)$, a contradiction since $X$ is extreme. Thus $\Delta_2 \nsubseteq X_1$, and this proves our claim that $\Delta_1 \subseteq X_1$. Similarly there is a restraint $X_2$ with $\Delta_2 \subseteq X_2$; and so $X_1 \cup X_2 \cup X = \Sigma$, contrary to (8.4). Thus our supposition that $W_1, W_2$ both have length at least that of $P$ was false, and the result follows. $\square$

\stmt
(8.8) \textit{If $X_1, X_2$ are extreme restraints and $X_1 \cap X_2 \ne \emptyset$ then $X_1 \cap X_2$ and $X_1 \cup X_2$ are simply-connected.}

\textit{Proof.} We shall first prove that $X_1 \cap X_2$ is arc-wise connected. Let its arc-wise connected components be $D_1, \dots, D_k$. We may assume that

\hspace*{12pt}(1) $D_h \cap bd(X_1) \ne \emptyset$, for $1 \le h \le k$.

\textit{Subproof.} We may assume that $D_h \ne X_2$, for otherwise $X_2 \subseteq X_1$, $k = 1$ and we are finished. Choose $a \in D_h$ and $b \in X_2 - D_h$, and let $\phi$ be a curve with $\phi(0) = a$, $\phi(1) = b$ and $|\phi| \subseteq X_2$. Choose $\alpha$ with $0 \le \alpha \le 1$ maximal so that $\phi(\beta) \in X_1$ for $0 \le \beta \le \alpha$ (this is possible since $X_1$ is closed). Then $\alpha < 1$, and so $\phi(\alpha) \in bd(X_1)$; and moreover, $\phi(\beta) \in X_1 \cap X_2$ for $0 \le \beta \le \alpha$, and so $\phi(\alpha) \in D_h$. This proves (1).
